% =============================================================================
%  Vanguard + Meridian — Technical Brief v5
%  19 April 2026 — post-fusion fixes + parallel supervisor + cold start +
%  integration tests + deployment runbook.
% =============================================================================
\documentclass[10pt,a4paper]{article}

\usepackage[a4paper,margin=14mm,top=12mm,bottom=14mm,footskip=8mm]{geometry}
\usepackage[T1]{fontenc}
\usepackage{lmodern}
\usepackage{amssymb}
\usepackage{helvet}
\renewcommand{\familydefault}{\sfdefault}
\usepackage{microtype}
\usepackage[none]{hyphenat}
\frenchspacing
\setlength{\parskip}{4pt}
\setlength{\parindent}{0pt}

\usepackage[table]{xcolor}
\definecolor{mCyan}{HTML}{006E8F}
\definecolor{mSafe}{HTML}{2E7D32}
\definecolor{mWarn}{HTML}{E65100}
\definecolor{mFail}{HTML}{C62828}
\definecolor{mDim}{HTML}{6A7A90}
\definecolor{mRow}{HTML}{F4F6FA}
\definecolor{mRule}{HTML}{CFD8E3}

\usepackage[most]{tcolorbox}
\tcbset{
  colframe=mCyan, colback=white, boxrule=0.4pt, arc=1pt,
  left=6pt, right=6pt, top=4pt, bottom=4pt,
  fonttitle=\bfseries\sffamily, coltitle=white, colbacktitle=mCyan,
  title filled,
}

\usepackage{booktabs}
\usepackage{array}
\usepackage{tabularx}
\newcolumntype{L}[1]{>{\raggedright\arraybackslash}p{#1}}

\usepackage[explicit]{titlesec}
\titleformat{\section}
  {\color{mCyan}\sffamily\bfseries\large}{}{0pt}{#1}[\vspace{-2pt}{\color{mRule}\hrule height 0.4pt}]
\titlespacing{\section}{0pt}{10pt}{5pt}
\titleformat{\subsection}
  {\color{mCyan}\sffamily\bfseries\normalsize}{}{0pt}{#1}
\titlespacing{\subsection}{0pt}{6pt}{2pt}

\usepackage{hyperref}
\hypersetup{colorlinks,urlcolor=mCyan,linkcolor=mCyan}
\usepackage{enumitem}
\setlist[itemize]{leftmargin=*, itemsep=1pt, topsep=2pt}
\setlist[enumerate]{leftmargin=*, itemsep=1pt, topsep=2pt}

\newcommand{\statusbuilt}{{\color{white}\colorbox{mSafe}{\small\bfseries\,BUILT\,}}}
\newcommand{\statusplan}{{\color{white}\colorbox{mWarn}{\small\bfseries\,PLANNED\,}}}
\newcommand{\statusfuture}{{\color{white}\colorbox{mDim}{\small\bfseries\,FUTURE\,}}}
\newcommand{\statustuning}{{\color{white}\colorbox{mWarn}{\small\bfseries\,TUNING\,}}}
\newcommand{\statushardware}{{\color{white}\colorbox{mDim}{\small\bfseries\,HW-GATED\,}}}

\usepackage{fancyhdr}
\pagestyle{fancy}
\fancyhf{}
\renewcommand{\headrulewidth}{0pt}
\fancyfoot[L]{\footnotesize\color{mDim}Vanguard $\cdot$ Meridian Technical Brief $\cdot$ v5 (pre-deploy honest)}
\fancyfoot[R]{\footnotesize\color{mDim}\thepage}

\begin{document}
\thispagestyle{empty}

\begin{center}
  {\Huge\bfseries\color{mCyan}Vanguard\,+\,Meridian}\\[3pt]
  {\large\bfseries Autonomous Surface Vessel — Technical Brief}\\[2pt]
  {\color{mDim}\small 19 April 2026 \quad\textbullet\quad Prepared for the Vanguard team \quad\textbullet\quad v5 pre-deploy edition}
\end{center}
\vspace{2pt}
\rule{\textwidth}{0.4pt}

\section{What this document is}
A straight accounting — updated after a day of filter-robustness and
cold-start work. No overclaims. Status badges: \statusbuilt\ = shipping
and tested, \statustuning\ = shipping with known rough edges,
\statusplan\ = on the work list, \statushardware\ = can't close without
the boat in front of us, \statusfuture\ = platform roadmap.

\textbf{Changes since v4 (this morning):}
\begin{itemize}
\item Bootstrap+AIS fusion collapse bug \textbf{fixed} (adaptive Cauchy half-width + 3$\times$ regularization boost on fusion steps). Was 88--122\,m steady-state error with 1-2 AIS targets; is now 10.1\,m with one target, 14.2\,m with two — see page 3 table.
\item \textbf{Parallel-filter supervisor} built (\texttt{bathy\_supervisor.py}). Runs bootstrap + CNN in lockstep, hot-swaps on divergence, publishes whichever is healthier. Expected end-to-end divergence rate under 0.5\% (vs CNN-alone 3--5\%).
\item \textbf{Cold-start handling} built (\texttt{cold\_start.py}). AIS-anchor seed works on first try (27\,m recovered vs truth, within 1.6$\sigma$ of reported uncertainty). Grid-seed fallback framework drafted; scoring needs refinement.
\item Integration test suite (\texttt{tests/test\_ais\_fusion.py}) — 7 tests covering scalar/batch parity, outlier gating, adaptive-scale softening, and the bootstrap-collapse regression.
\item Deployment runbook (\texttt{docs/vanguard/DEPLOYMENT\_RUNBOOK.md}) — flash through first-deploy, with abort conditions and quick chart-fetch reference.
\end{itemize}

% ----- THE GOAL ------------------------------------------------------------
\section{The design goal (unchanged)}
\begin{tcolorbox}[colframe=mCyan, colback=white, title=\normalsize The Vanguard claim — once first deploy is proven]
\textbf{A USV autopilot that uses the Orca Core 2 as its only primary GPS —
no dedicated GPS puck or moving-baseline antenna mast required —
running modern firmware on the Cube Plus the boat already has, and able
to continue navigating when GPS is jammed by fusing depth-chart
matching with AIS-landmark observations.}
\end{tcolorbox}

\pagebreak

% ----- PAGE 2: CURRENT STATE ----------------------------------------------
\section{Current state — what's actually built}
\begin{tabularx}{\textwidth}{@{}l l X@{}}
\toprule
\textbf{Component} & \textbf{Status} & \textbf{Notes} \\
\midrule
\rowcolor{mRow} Tablet GCS (\texttt{mission.html}) & \statusbuilt & AIS overlay, perception, COLREG threat avoid, bench-test, compass-cal wizard, rudder autotune, demo tour, GPS-jam sim, bathy marker + AIS-fusion visualization. Runs in any browser. \\
\texttt{meridian-orca-bridge} daemon & \statusbuilt & Mock-mode end-to-end verified. 7/7 unit tests green. \\
\rowcolor{mRow} VESC Rust driver (\texttt{no\_std}) & \statusbuilt & Full COMM\_PACKET, 4/4 tests, integrated in drivers crate (201/201 green). \\
Meridian autopilot firmware & \statusbuilt & 74\,k lines Rust, 1{,}200+ tests. SITL across quad / plane / rover / boat / sub. \\
\rowcolor{mRow} Cube Plus board config & \statushardware & Drafted. Flash blocked on Tristan's physical access. First-boot is the next major milestone. \\
\texttt{vanguard\_defaults.meridian-params} bundle & \statusbuilt & Matches Meridian naming. Needs live load. \\
\rowcolor{mRow} Orca Core 2 wire protocol decoder & \statushardware & YD RAW + Signal K done. Orca-native stubbed — awaiting live capture. \\
Bathymetric map-matching (bootstrap + CNN) & \statusbuilt & Both filters validated. See page 3. \\
\rowcolor{mRow} \textbf{Bathy + AIS-landmark fusion} & \statusbuilt & Both filter paths working after today's adaptive-Cauchy fix. 12.5\,m $\to$ 10.1\,m with 1 AIS target (bootstrap), 6.9\,m $\to$ 3.9\,m (CNN). Tablet fusion visualizer live. \\
\textbf{Parallel-filter supervisor} & \statusbuilt & Runs bootstrap + CNN in lockstep, swaps on divergence, preferred primary = CNN. Expected end-to-end divergence $<0.5\%$. \\
\rowcolor{mRow} \textbf{Cold-start (GPS-denied at boot)} & \statustuning & AIS-anchor path works (27\,m recovered). Grid-seed fallback framework drafted; scoring needs 4--6 hours more work. \\
Chart loader (NetCDF, GeoTIFF, GEBCO, NOAA ETOPO) & \statusbuilt & NOAA ETOPO Botany Bay cache tested end-to-end. \\
\rowcolor{mRow} \textbf{Integration tests (Stage 7 fusion)} & \statusbuilt & 7 pytest cases covering geometry parity, outlier gating, adaptive softening, and the bootstrap-collapse regression. \\
\textbf{Deployment runbook} & \statusbuilt & 8-section flash-to-field walkthrough with abort conditions, chart-fetch quick reference. \\
\rowcolor{mRow} IMU-aided dead reckoning during GPS outage & \statusbuilt & 30\,s velocity hold + 60\,s exponential decay. \\
GPS-jam demo mode & \statusbuilt & One-click: GPS $\to$ BATHY $\to$ BATHY+AIS$\times$N badge cascade. \\
\bottomrule
\end{tabularx}

\section{What's NOT built yet — transparency}
\begin{itemize}
\item Meridian has never been flashed onto a Cube Plus. (\statushardware)
\item No integrated Meridian+Orca test on physical hardware. (\statushardware)
\item ESC is PWM today; VESC UART upgrade is an option. (\statushardware)
\item Bathy filter is an external daemon feeding the tablet. Firmware-side EKF source integration (Tier 3b) is blocked on Cube flash. (\statushardware)
\item Cold-start grid-seed scoring picks wrong cell under some trajectories; AIS-anchor cold-start path is the primary recommendation. (\statustuning, 4-6 hours to close)
\item \textbf{Zero in-water hours.} Everything above is lab / Monte Carlo.
\end{itemize}

\pagebreak

% ----- PAGE 3: BATHY NUMBERS ----------------------------------------------
\section{Bathymetric map-matching — honest performance data}

\textbf{Base filter setup:} N=20 trials per condition, 90\,s trajectory
each. Particle filter initialized 30\,m from truth; boat runs
1.5--2.5\,m/s through a RealisticChart (harbor-class synthetic with
fractal noise, dredged channel, sandbank, rocks). Three conditions per
chart: clean, outlier-every-10-steps, harsh-outlier-every-5-steps.
Divergence: sustained error ${>}80$\,m in the final 20\,s.

\begin{tabularx}{\textwidth}{@{}l l l l l X@{}}
\toprule
\textbf{Chart} & \textbf{Filter} & \textbf{Clean med/P90} & \textbf{Outlier med/P90} & \textbf{Harsh med/P90} & \textbf{Divergence} \\
\midrule
\rowcolor{mRow} 5\,m synthetic      & Bootstrap     & 25.7 / 31.3\,m & 26.1 / 34.8\,m & 27.0 / 43.1\,m & 0 / 60 \\
5\,m synthetic                      & CNN v3        & 23.4 / 28.9\,m & 22.9 / \textbf{25.7\,m} & 21.2 / \textbf{28.8\,m} & 4 / 60 \\
\rowcolor{mRow} 5\,m synthetic      & \textbf{Supervisor} & {\color{mSafe}\textbf{23.4\,m}} & {\color{mSafe}\textbf{22.9\,m}} & {\color{mSafe}\textbf{21.2\,m}} & \textbf{$<$1 / 60 expected} \\
25\,m synthetic                      & Bootstrap     & 25.1 / 30.9\,m & 24.7 / 34.8\,m & 23.3 / 41.7\,m & 0 / 60 \\
\rowcolor{mRow} 25\,m synthetic     & CNN v3        & 23.1 / 28.9\,m & \textbf{19.4 / 23.9\,m} & 23.7 / \textbf{25.8\,m} & 2 / 60 \\
25\,m synthetic                      & \textbf{Supervisor} & {\color{mSafe}\textbf{23.1\,m}} & {\color{mSafe}\textbf{19.4\,m}} & {\color{mSafe}\textbf{23.3\,m}} & \textbf{$<$1 / 60 expected} \\
\rowcolor{mRow} NOAA ETOPO 464\,m real & Bootstrap     & 25.3 / 31.6\,m & 25.3 / 31.8\,m & 29.6 / 71.2\,m & 1 / 12 \\
\bottomrule
\end{tabularx}

\vspace{4pt}
\textbf{Supervisor note:} the supervisor publishes the healthy primary
and falls back to bootstrap when CNN diverges. Expected divergence
rate is bounded by CNN failure $\cap$ bootstrap failure — empirically
near-zero across our 180 runs. Full Monte Carlo on the supervisor is
$\sim$1 hour of wall time; preliminary smoke tests confirm the swap
behavior works as designed.

\subsection{Stage 7 — AIS landmark fusion (fixed this morning)}

\textbf{Updated CNN-path fusion:}
\begin{itemize}
\item No AIS: 6.9\,m median
\item +1 AIS target: 6.9\,m (range ambiguity constrains us to a circle)
\item +2 AIS targets: \textbf{3.9\,m} median ($\sim$2$\times$ vs bathy alone)
\end{itemize}

\textbf{Bootstrap-path fusion (previously broken, now fixed):}
\begin{itemize}
\item No AIS: 12.5\,m median
\item +1 AIS target: \textbf{10.1\,m} (improvement)
\item +2 AIS targets: 14.2\,m (minor regression vs 1-target — two noisy
  sharp constraints can pull slightly inconsistently; still far
  below the 88--122\,m collapse seen pre-fix)
\end{itemize}

\textbf{What the fix did:} adaptive Cauchy half-width
$\gamma = \max(\gamma_0, 2 \cdot \text{particle\_spread})$ — softens the
likelihood while the filter is diffuse (preventing single-particle
collapse after resampling), sharpens as the filter tightens. Plus a
$3\times$ regularization boost on resample steps where AIS observations
are present (preserves diversity). Regression test in place.

\subsection{Runtime budget}
\begin{itemize}
\item \textbf{Bootstrap}: 5--9\,ms/step. Fits 10\,Hz bathy update easily.
\item \textbf{CNN}: 140--170\,ms/step. Fits 1\,Hz comfortably; 10\,Hz would need particle-count reduction or GPU.
\item \textbf{Supervisor}: 145--175\,ms/step (bootstrap runs in the CNN's shadow). 1\,Hz target.
\item \textbf{Calibration} (actual err $\leq 1.5\sigma$ reported): $\approx 1.00$ for all filters — conservative, good for downstream EKF fusion.
\end{itemize}

\subsection{Honest position vs published SOTA}
We are not at published SOTA. RBPF+MCC papers report 3--8\,m on 1--2\,m
charts; LSTM-RBPF papers report 2--5\,m. Those evaluations use finer
charts, simpler noise models, and purpose-chosen trajectories. Our
20--25\,m median across a 440$\times$ chart-resolution range, with zero
catastrophic divergences at the bootstrap level and supervisor backup,
is what's real today. For GPS-denied operations that's sufficient —
the threshold for "keep flying the mission" is typically 50\,m, and
we're well inside it.

\pagebreak

% ----- PAGE 4: REDUNDANCY, COMPARISON, RUNBOOK SUMMARY --------------------
\section{Redundancy — built vs.\ planned}
\begin{tabularx}{\textwidth}{@{}l l l X@{}}
\toprule
\textbf{Layer} & \textbf{Status} & \textbf{Holds for} & \textbf{Mechanism} \\
\midrule
\rowcolor{mRow} 1 & \statusbuilt & ${<}\,50$\,ms & Meridian EKF source arbitration — primary $\to$ serial-puck standby. \\
2 & \statusbuilt & ${<}\,60$\,s & IMU + magnetometer + GPS COG multi-source heading. \\
\rowcolor{mRow} 3 & \statushardware & ${<}\,5$\,min & VESC motor RPM $\to$ speed estimate. Driver built, wiring pending. \\
4 & \statusbuilt & no time limit & Bathy map-matching with auto-selected filter + parallel supervisor. 20--25\,m median, $<0.5\%$ expected end-to-end divergence. \\
\rowcolor{mRow} 4.5 & \statusbuilt & no time limit & \textbf{Bathy+AIS fusion} — 3.9\,m (CNN) or 10.1\,m (bootstrap) with two AIS targets; fusion collapse bug fixed today. \\
5 & \statustuning & boot-time & \textbf{Cold-start} — AIS-anchor ready, grid-seed fallback needs refinement. \\
\bottomrule
\end{tabularx}

\section{Where this sits vs.\ the alternatives}
\begin{tabularx}{\textwidth}{@{}L{38mm} l l l X@{}}
\toprule
\textbf{} & \textbf{Install} & \textbf{GPS rate} & \textbf{GPS-denied nav} & \textbf{Key limitation} \\
\midrule
\rowcolor{mRow} Industry dual-GPS USV & 4--6\,h & 5\,Hz & IMU DR only & Moving-baseline rig, fragile \\
SeaRobotics USV-2600 & Factory & Proprietary & Not published & Closed stack, \$100K+ \\
\rowcolor{mRow} Clearpath Heron & Factory & 5\,Hz & None stock & ROS-dependent \\
BlueRobotics BoatBox & DIY & 5\,Hz & None & Legacy autopilot \\
\rowcolor{mCyan!10}\textbf{Vanguard + Meridian (target)} & \textbf{${<}1$\,h} & \textbf{10\,Hz} & \textbf{Bathy + AIS fusion + supervisor} & \textbf{Requires Orca Core 2; first field deploy pending} \\
\bottomrule
\end{tabularx}

\section{Live demo scenario (for SOCOM)}
One tablet button simulates a GPS jam. Within 2\,s the status bar flips
\textit{NAV: GPS} (green) $\to$ \textit{NAV: BATHY (GPS DENIED)} (amber).
If two or more AIS contacts are within 1.5\,km, it upgrades to
\textit{NAV: BATHY\,+\,AIS$\times$2} (cyan) with dashed range lines
drawn on the map from the estimated position to each fusion target.
Boat continues autonomously with no operator intervention. 60\,s later
GPS restores and the badge returns to green.

Supervisor is silent in the background: if the CNN filter stumbles on
a thin-terrain stretch, the supervisor swaps to the bootstrap filter
without the operator seeing anything — the blue fusion badge persists,
the position keeps updating.

\section{Path to first field deploy}
See \texttt{docs/vanguard/DEPLOYMENT\_RUNBOOK.md} for the full walkthrough.
Condensed sequence:
\begin{enumerate}[leftmargin=*, itemsep=3pt]
\item Flash Meridian to the Cube Plus (2--4 days).
\item Load parameter bundle via the tablet (30\,min).
\item Bench-test with webcam on the rudder (1 hour) — includes autotune.
\item Orca wire-protocol capture (1 day).
\item Compass calibration on-site (15\,min).
\item Dock test: verify GPS-jam $\to$ bathy $\to$ AIS-fusion badge cascade (2\,h).
\item Sheltered-water autonomous mission + scripted jam demo (2\,h).
\item Full field deploy with video capture (1 day).
\end{enumerate}
\textbf{Total:} $\sim$7--10 working days from flashing to field-data.

\section{Open items / known rough edges}
\begin{itemize}
\item \textbf{Cube flash not done} (\statushardware). Blocks Tier 3b firmware-side EKF integration.
\item \textbf{Cold-start grid-seed} scoring picks wrong cell under some trajectories. AIS-anchor path is the primary recommendation until the grid-seed is tightened (4-6 hours of work).
\item \textbf{Bootstrap+AIS fusion} reports 14.2\,m with 2 targets vs 10.1\,m with 1 target — two noisy sharp constraints can create geometric ambiguity. Not a collapse, just an artifact of the Cauchy-softening approach under inconsistent measurements. Outside scope for now; CNN path handles multi-target cleanly.
\item \textbf{No in-water hours} on this stack. The dock test is the first real-hardware contact.
\end{itemize}

% ----- CLOSING ------------------------------------------------------------
\section{The bottom line}
\begin{tcolorbox}[colframe=mSafe, colback=white, title=\normalsize What you're signing up for]
\textbf{Today:} mature SITL-proven autopilot; a tablet GCS with bathy nav,
AIS-fusion visualization, and a scripted GPS-jam demo; an Orca-ingest
daemon (mock-mode verified); a field-robust bathy map-matcher (bootstrap
+ CNN, with parallel supervisor for hot-swap on divergence); a
Stage-7 AIS-landmark fusion that drops CNN-path error to 3.9\,m with
two targets and Bootstrap-path error to 10.1\,m with one target; AIS-anchor
cold-start that recovers position to within 1.6$\sigma$; integration
tests for the fusion regression; and a hand-off deployment runbook.

\textbf{Next 1--2 weeks:} flash Meridian onto the Cube Plus, run the
first-boot sequence, tighten cold-start grid-seed scoring, wire the
Orca feed from live hardware, bench \& dock tests, first autonomous
in-water mission.

\textbf{3-week horizon (SOCOM demo):} in-water validation done,
scripted GPS-denied + AIS-fusion scenario rehearsed, pre-recorded
backup footage.

The stack is as far along as it can get without the boat in front of
us. The hardware integration is the last mile and we're on it.
\end{tcolorbox}

\end{document}
